\section{Geometric structure of class regions in matrix space}
\label{app:geometric}

This appendix preserves a third theoretical result from the
construction of \citet{armenta2021representation} that is not
load-bearing for either Study~1 (Theorem~\ref{thm:maxinv}) or
Study~2 (Theorem~\ref{thm:pyth}). We state it here together with
its geometric implications and two candidate experimental directions
that would promote it to a third theoretical leg of the
canonical-representation argument in a future revision.

\paragraph{Class regions in matrix space.}
For a network $(W, f)$ on a $C$-class classification task, define the
\emph{class region} for class $j$ as
\[
\mathcal{M}_j \;=\; \Bigl\{ M \in \mathrm{Mat}_\R(C, d+1)
   \;:\; (M \cdot \mathbf{1}_{d+1})_j > (M \cdot \mathbf{1}_{d+1})_i
   \text{ for all } i \neq j \Bigr\}.
\]
By the completeness invariant, $\Mfx \in \mathcal{M}_j$ if and only
if the network classifies $x$ as class $j$. The set $\mathcal{M}_0
= \mathrm{Mat}_\R(C, d+1) \setminus \bigcup_{j=1}^{C} \mathcal{M}_j$
of indeterminate matrices is of measure zero
(Theorem~4.3 of \citealt{armenta2021representation}), so we may
restrict attention to the $\mathcal{M}_j$ partition.

\begin{theorem}[Convexity of class regions, Theorem~4.4 of
\citealt{armenta2021representation}]
\label{thm:convexity}
For each $j > 0$, the class region $\mathcal{M}_j$ is convex in
$\mathrm{Mat}_\R(C, d+1)$.
\end{theorem}

\begin{proof}
For matrices $A, B \in \mathcal{M}_j$ and any $\lambda \in [0, 1]$,
the row-sum identity gives $((1-\lambda) A + \lambda B) \cdot
\mathbf{1}_{d+1} = (1-\lambda)(A \mathbf{1}_{d+1}) +
\lambda(B \mathbf{1}_{d+1})$. The argmax of a non-negative convex
combination of two vectors that share a common argmax is that same
argmax, so $((1-\lambda) A + \lambda B) \in \mathcal{M}_j$.
\end{proof}

\paragraph{Geometric implications.}
Theorem~\ref{thm:convexity} says the matrix-space partition into
class regions is well-behaved enough for classical convex-set tools
(centroid, support, projection) to apply directly to knowledge
matrices, with no further structural assumptions. We do not exercise
this property in the studies of this paper.

\paragraph{Toward a third theoretical leg.}
Theorem~\ref{thm:convexity} could be promoted to a third main theorem
in a future revision of this paper, supporting a third experimental
leg of the canonical-representation argument. We outline two natural
candidate directions and leave both as future work:

\begin{enumerate}
\item \emph{$k$-nearest-neighbour classification in matrix space.}
The convex class region admits a clean canonical centroid, against
which $k$-NN-on-knowledge-matrices is well-defined without any
representation alignment step. The corresponding $k$-NN-on-penultimate
baseline requires CKA-style alignment to compare across networks
(\citealt{kornblith2019similarity}), which the matrix-space version
sidesteps by Theorem~\ref{thm:maxinv}.

\item \emph{Convex-hull federated aggregation.}
For clients training networks within a function-equivalence class,
averaging the per-client knowledge matrices produces a result that
lies inside the corresponding class region $\mathcal{M}_j$ by
convexity (Theorem~\ref{thm:convexity}). Penultimate-feature
averaging cannot give the analogous guarantee, because penultimate
features are not invariants of the function
(Theorem~\ref{thm:maxinv}). A federated-learning experiment in
this direction would couple Theorem~\ref{thm:maxinv} (Study~1)
with Theorem~\ref{thm:convexity} into a single demonstration.
\end{enumerate}
